\newpage
\phantomsection
\label{prf:uniqueness-of-the-maximum}

\begin{remark*}[Return]
\hyperref[prop:uniqueness-of-the-maximum]{Return to Theorem}
\end{remark*}

\begin{theorem*}[Uniqueness of the Maximum]
If a set $A \subseteq \mathbb{R}$ has a maximum, then that maximum is unique.
\end{theorem*}

\begin{proof}
\textbf{Professional Standard Proof.}~\\
Let $m_1, m_2 \in A$ be maxima of $A$. By the definition of maximum, $m_1$ is an upper bound of $A$; since $m_2 \in A$, it follows that $m_2 \le m_1$. Similarly, $m_2$ is an upper bound of $A$; since $m_1 \in A$, it follows that $m_1 \le m_2$. By the antisymmetry of the $\le$ relation on $\mathbb{R}$, $m_1 \le m_2$ and $m_2 \le m_1$ imply $m_1 = m_2$.
\end{proof}

\begin{proof}
\textbf{Detailed Learning Proof.}~\\

\textbf{Step 1.} Assume the existence of two maxima.
Let $A \subseteq \mathbb{R}$ be a non-empty set. Suppose that $m_1$ and $m_2$ are both maxima of $A$. By the definition of maximum, this implies:
\[
m_1 \in A \quad \text{and} \quad m_2 \in A.
\]

\textbf{Step 2.} Use the upper bound property of the first maximum.
Since $m_1$ is a maximum, it is an upper bound for $A$. Therefore, every element $a \in A$ satisfies $a \le m_1$. Because $m_2$ is an element of $A$, it follows that:
\[
m_2 \le m_1.
\]

\textbf{Step 3.} Use the upper bound property of the second maximum.
Conversely, since $m_2$ is a maximum, every element $a \in A$ satisfies $a \le m_2$. Because $m_1$ is an element of $A$, it follows that:
\[
m_1 \le m_2.
\]

\textbf{Step 4.} Apply antisymmetry to conclude uniqueness.
The results of Step 2 and Step 3 provide the pair of inequalities:
\[
m_1 \le m_2 \quad \text{and} \quad m_2 \le m_1.
\]
By the antisymmetry of the real ordering, this forces $m_1 = m_2$. Thus the maximum is unique.
\end{proof}

\begin{remark*}[Proof structure]
This is a uniqueness proof by antisymmetry. It assumes two candidates for the maximum and uses the double-inclusion logic of upper bounds to force their equality.
\end{remark*}

\begin{remark*}[Dependencies]\hfill
\begin{itemize}
  \item \hyperref[def:maximum]{Definition of Maximum}
  \item \hyperref[ax:antisymmetry]{Antisymmetry of $\mathbb{R}$}
\end{itemize}
\end{remark*}